\subsection{Reducci\'on de dimensionalidad}
Se analizaron 1906 observaciones de entrenamiento y 477 observaciones de prueba, descritas por 64 coeficientes MFCC continuos. Antes de aplicar los m\'etodos de reducci\'on de dimensionalidad, se estandarizaron las variables \texttt{mel\_0}--\texttt{mel\_63} mediante media cero y varianza unitaria. Esta etapa fue requerida debido a que PCA, t-SNE y UMAP son sensibles a la escala de entrada.

Se ajust\'o PCA completo sobre el conjunto de entrenamiento con el objetivo de estimar la varianza acumulada por componente principal. El umbral de 95\% de varianza acumulada fue alcanzado con 11 componentes. Posteriormente, se calcul\'o una proyecci\'on PCA en dos dimensiones para compararla contra una t\'ecnica de variedades no lineal. En esta ejecuci\'on, la comparaci\'on se realiz\'o con t-SNE 2D, siguiendo una inicializaci\'on basada en PCA cuando el algoritmo lo permiti\'o.

\begin{table}[H]
\centering
\caption{Comparaci\'on cuantitativa de m\'etodos de reducci\'on de dimensionalidad.}
\begin{tabular}{lrrrr}
\hline
M\'etodo & Tiempo (s) & Varianza 2D & Trustworthiness@10 & Corr. distancias \\
\hline
PCA 2D & 0.002 & 0.671 & 0.867 & 0.947 \\
t-SNE 2D & 13.948 & N/A & 0.987 & 0.590 \\
\hline
\end{tabular}
\end{table}

Se observ\'o que PCA proporcion\'o una referencia lineal interpretable mediante varianza explicada y error de reconstrucci\'on, mientras que la t\'ecnica de variedades fue evaluada mediante preservaci\'on local y correlaci\'on de distancias, ya que no dispone de una raz\'on de varianza explicada equivalente. La separaci\'on visual por \texttt{species\_id} fue inspeccionada en el plano bidimensional para identificar posibles agrupamientos naturales entre las cinco clases.
